\documentclass[9pt,a4paper,landscape]{extarticle}
\usepackage[T1]{fontenc}
\usepackage[utf8]{inputenc}
\usepackage[french]{babel}
\usepackage{lmodern,geometry,array,longtable,graphicx,multirow,amssymb}
\usepackage[table]{xcolor}
\usepackage{hyperref,fancyhdr}
\geometry{margin=10mm,bottom=12mm}
\hypersetup{colorlinks=true,urlcolor=blue,linkcolor=black}
\definecolor{ink}{HTML}{163245}
\definecolor{pale}{HTML}{E9F0F4}
\pagestyle{fancy}\fancyhf{}\fancyfoot[L]{\scriptsize Recherche citoyenne - hypothèses M0.1 - 26 octobre 2026}
\fancyfoot[R]{\scriptsize\thepage}\renewcommand{\headrulewidth}{0pt}
\setlength{\parindent}{0pt}\setlength{\parskip}{4pt}
\newcommand{\rot}[1]{\rotatebox{90}{\hspace{1mm}\parbox{39mm}{\raggedright\fontsize{8}{9}\selectfont #1}\hspace{1mm}}}
\newcommand{\rel}{\ensuremath{\circ^{?}}}
\newcommand{\proc}{\ensuremath{\triangle^{?}}}
\newcommand{\nd}{\textcolor{gray}{ND}}
\begin{document}
{\Large\bfseries\color{ink} Matrice des compétences - M0.1}\hfill{\large Scénario Codex}\par
\textbf{Projet de formation syndicale : hypothèse à examiner.} Adaptation de la forme du guide, § 2.4 et annexe III.\par
{\small Préparation Codex; base documentaire commune. Gemini réinterprète le travail Codex. Aucun accord entre IA ne vaut validation indépendante.}\par
\vspace{1mm}\begingroup\setlength{\tabcolsep}{1mm}\renewcommand{\arraystretch}{1.05}\fontsize{8}{9}\selectfont
\begin{tabular}{|p{83mm}|>{\centering\arraybackslash}m{7mm}|>{\centering\arraybackslash}m{7mm}|>{\centering\arraybackslash}m{8mm}||>{\centering\arraybackslash}m{11mm}|>{\centering\arraybackslash}m{11mm}|>{\centering\arraybackslash}m{11mm}|>{\centering\arraybackslash}m{11mm}|>{\centering\arraybackslash}m{11mm}|>{\centering\arraybackslash}m{11mm}|>{\centering\arraybackslash}m{11mm}|>{\centering\arraybackslash}m{11mm}|>{\centering\arraybackslash}m{11mm}|>{\centering\arraybackslash}m{11mm}|>{\centering\arraybackslash}m{11mm}|}
\hline
\rowcolor{ink}\multicolumn{15}{|c|}{\color{white}\bfseries MATRICE DES COMPÉTENCES - RELATIONS PROPOSÉES}\tabularnewline\hline
\multicolumn{4}{|c||}{\textbf{COMPÉTENCES PARTICULIÈRES}} & \multicolumn{6}{c|}{\textbf{COMPÉTENCES GÉNÉRALES}} & \multicolumn{5}{c|}{\textbf{PROCESSUS HYPOTHÉTIQUE}}\tabularnewline\hline
\textbf{Mandat élu syndical}\par Service ou résultat proposé en ligne.\par Les codes renvoient à la table de correspondance. & \rot{Numéro proposé} & \rot{Type d’objectif} & \rot{Durée (h)} & \rot{Analyser l’information syndicale} & \rot{Communiquer l’information syndicale} & \rot{Collaborer avec les acteurs du mandat} & \rot{Délimiter son intervention syndicale} & \rot{Développer ses apprentissages du mandat} & \rot{Mobiliser les ressources d’un réseau syndical} & \rot{Recueillir les préoccupations des membres} & \rot{Analyser les enjeux du mandat} & \rot{Préparer l’intervention syndicale} & \rot{Réaliser l’intervention syndicale} & \rot{Examiner les suites de l’intervention}\tabularnewline\hline
\multicolumn{4}{|r||}{Code / numéro} & CG01/1 & CG02/2 & CG03/3 & CG04/4 & CG05/5 & CG06/6 & P1 & P2 & P3 & P4 & P5\tabularnewline\hline
\multicolumn{4}{|r||}{Type d’objectif proposé} & C? & C? & C? & C? & S? & C? & \multicolumn{5}{c|}{Sans objet pour les étapes}\tabularnewline\hline
\multicolumn{4}{|r||}{Durée (h)} & \nd & \nd & \nd & \nd & \nd & \nd & \multicolumn{5}{c|}{Sans durée additionnelle}\tabularnewline\hline
\rowcolor{pale}
\rule{0pt}{5.2mm}\textbf{CP-C01} Préparer l’accueil des membres & 7 & C? & \nd & \rel & \rel & \rel & \rel &  &  & \proc & \proc & \proc & \proc & \proc\tabularnewline\hline
\rule{0pt}{5.2mm}\textbf{CP-C02} Réaliser une consultation des membres & 8 & C? & \nd & \rel & \rel & \rel & \rel &  &  & \proc & \proc & \proc & \proc & \proc\tabularnewline\hline
\rowcolor{pale}
\rule{0pt}{5.2mm}\textbf{CP-C03} Préparer une instance syndicale & 9 & C? & \nd & \rel & \rel & \rel & \rel &  &  &  & \proc & \proc & \proc & \proc\tabularnewline\hline
\rule{0pt}{5.2mm}\textbf{CP-C04} Traiter un dossier de représentation & 10 & C? & \nd & \rel & \rel & \rel & \rel & \rel & \rel & \proc & \proc & \proc & \proc & \proc\tabularnewline\hline
\rowcolor{pale}
\rule{0pt}{5.2mm}\textbf{CP-C05} Accompagner un membre en difficulté & 11 & C? & \nd & \rel & \rel & \rel & \rel &  & \rel & \proc & \proc & \proc & \proc & \proc\tabularnewline\hline
\rule{0pt}{5.2mm}\textbf{CP-C06} Traiter un différend syndical & 12 & C? & \nd & \rel & \rel & \rel & \rel &  &  & \proc & \proc & \proc & \proc & \proc\tabularnewline\hline
\rowcolor{pale}
\rule{0pt}{5.2mm}\textbf{CP-C07} Coordonner le travail de l’équipe syndicale & 13 & C? & \nd & \rel & \rel & \rel & \rel &  & \rel & \proc & \proc & \proc & \proc & \proc\tabularnewline\hline
\rule{0pt}{5.2mm}\textbf{CP-C08} Organiser une action collective & 14 & C? & \nd & \rel & \rel & \rel & \rel &  & \rel & \proc & \proc & \proc & \proc & \proc\tabularnewline\hline
\rowcolor{pale}
\rule{0pt}{5.2mm}\textbf{CP-C09} Préparer une prise de position syndicale & 15 & C? & \nd & \rel & \rel & \rel & \rel &  & \rel &  & \proc & \proc & \proc & \proc\tabularnewline\hline
\rule{0pt}{5.2mm}\textbf{CP-C10} Négocier les moyens de libération syndicale & 16 & C? & \nd & \rel & \rel & \rel & \rel &  &  &  & \proc & \proc & \proc & \proc\tabularnewline\hline
\rowcolor{pale}
\rule{0pt}{5.2mm}\textbf{CP-C11} Préparer une réponse à un projet de sous-traitance & 17 & C? & \nd & \rel & \rel & \rel & \rel & \rel & \rel &  & \proc & \proc & \proc & \proc\tabularnewline\hline
\rule{0pt}{5.2mm}\textbf{CP-C12} Organiser une action de solidarité communautaire & 18 & C? & \nd &  & \rel & \rel & \rel &  & \rel & \proc & \proc & \proc & \proc & \proc\tabularnewline\hline
\rowcolor{pale}
\rule{0pt}{5.2mm}\textbf{CP-C13} Préparer la transmission d’un mandat syndical & 19 & C? & \nd & \rel & \rel & \rel & \rel & \rel & \rel & \proc & \proc & \proc & \proc & \proc\tabularnewline\hline
\end{tabular}\endgroup\par\vspace{1mm}
{\footnotesize \textbf{Lecture M0.1 :} $\circ^{?}$ = lien général/particulier proposé; $\triangle^{?}$ = étape proposée. Case vide = lien non retenu à ce stade, sans preuve d’absence.\par \textbf{Choix encore ouverts :} C? = objectif de comportement proposé; S? = objectif de situation proposé; ND = durée à déterminer, jamais zéro. Aucun total d’heures calculable.\par Les symboles pleins $\bullet$ et $\blacktriangle$ sont réservés aux applications choisies dans les objectifs de formation; aucune n’est décidée ici. L’ordre proposé reste à examiner.}
\clearpage{\Large\bfseries\color{ink} Mode d’emploi et décisions de construction}\par
\textbf{Statut :} M0.1 est une adaptation méthodologique avant terrain, pas M1 après données brutes. Les M0 du 26 octobre restent conservées pour comparer les changements.\par
\textbf{Limite reconnue :} Les premières M0 étaient des listes structurées. Elles ne distinguaient pas les deux types de compétences et ne comportaient pas cette grille de relations. Le regroupement automatique d’une compétence par famille ne suffisait pas.\par
\textbf{Point de départ :} La pré-AST compte 47 constats, dont 44 activités et trois conditions. Elle ne remplace pas un rapport d’AST issu d’un atelier. Les activités, fonctions et critères doivent encore être confrontés au terrain.\par
\textbf{Transformation :} Chaque ligne particulière vise un service ou un résultat distinct. Les activités transférables sont proposées en colonnes générales. Les indications et conditions restent dans la table de correspondance plutôt que dans des titres trop longs.\par
\textbf{Ordre et numéros :} Les six compétences générales précèdent les lignes particulières dans la numérotation proposée. Les lignes vont de situations délimitées vers des interventions collectives ou spécialisées, puis la transmission. La difficulté réelle et l’ordre d’apprentissage ne sont pas mesurés.\par
\textbf{Objectifs et durées :} C? est proposé quand un résultat observable peut être décrit. S? est proposé pour le développement de ses apprentissages, qui demandera une démarche et des critères de participation. Ces types et toutes les durées restent à décider après analyse des besoins.\par
\textbf{Processus :} Les cinq étapes décrivent une possibilité de cheminement syndical. Elles ne constituent pas une procédure constatée commune à tous les postes. Les triangles n’imposent ni ordre rigide ni déroulement linéaire.\par
\textbf{Provenance :} Les deux matrices sont préparées par Codex. Le classement de Gemini est une réinterprétation du travail Codex. Les nouvelles distinctions, reformulations et relations M0.1 sont des propositions ajoutées par Codex à la demande de Steve.\par
\textbf{Décision :} Steve examine les compétences, les relations et les critères; les données brutes et les entretiens pourront les contredire, les diviser ou les supprimer. Aucune clôture de 1B ni compétence validée n’est annoncée.\par
\textbf{Buts de travail du projet :} B1 - soutenir une intervention adaptée au mandat; B2 - soutenir la prise de fonction et la participation; B3 - développer l’apprentissage et la transmission. Ces buts sont proposés pour la recherche citoyenne. Ils ne constituent pas l’adoption d’un programme ministériel.\par
\textbf{Référence de forme :} Ministère de l’Éducation du Québec (2004), \emph{Guide de conception et de production d’un programme d’études professionnelles}, sections 1.2 et 2.1 à 2.5, p. 6-16; annexes III-IV, p. 55-58.\par\href{https://www.inforoutefpt.org/ministere_docs/adminInfo/guideProgramme/GuideConceptionProduction_FP.pdf}{Consulter le guide source}. Un mandat élu nécessite une adaptation du cadre professionnel.\par
\clearpage{\Large\bfseries\color{ink} Table de correspondance - compétences générales}\par
\small Les appuis ci-dessous sont des passages documentaires; les compétences et leur portée transférable sont inférées. Les buts B1-B3 sont ceux de la page précédente.\par
\renewcommand{\arraystretch}{1.15}\begin{longtable}{|p{55mm}|p{95mm}|p{110mm}|}\hline\rowcolor{pale}\textbf{Énoncé / origine} & \textbf{Indications et buts} & \textbf{Déterminants et limites}\\\hline\endhead
CG01 - Analyser l’information syndicale\par M0-C01, M0-C03, M0-C13 & Distinguer les faits, les récits et les informations manquantes; comparer les éléments utiles au mandat. Buts : B1, B3. & E001 (S01, p. 301, 304); E005 (S01, p. 306-307); E029 (S04, p. 167-168) Portée générale et nécessité des liens à vérifier.\\\hline
CG02 - Communiquer l’information syndicale\par M0-C02, M0-C05, M0-C06, M0-C09 & Adapter le message au public; expliquer les choix; vérifier la compréhension. Buts : B1, B2. & E006 (S01, p. 307); E008 (S01, p. 308-309); E019 (S02, p. 81); E047 (S04, p. 145) Portée générale et nécessité des liens à vérifier.\\\hline
CG03 - Collaborer avec les acteurs du mandat\par M0-C03, M0-C07, M0-C08 & Clarifier les contributions; organiser un relais; discuter d’un désaccord de rôles. Buts : B1, B2. & E003 (S01, p. 304); E014 (S03, p. 215-216, 219-220); E042 (S03, p. 229-232); E043 (S03, p. 230, 235) Portée générale et nécessité des liens à vérifier.\\\hline
CG04 - Délimiter son intervention syndicale\par M0-C01, M0-C07, M0-C08, M0-C11 & Distinguer son mandat de celui d’un autre acteur; choisir quand demander un avis ou passer le relais. Buts : B1, B2. & E007 (S01, p. 308); E009 (S01, p. 309-310); E016 (S03, p. 217); E041 (S03, p. 214, 221-222, 344) Portée générale et nécessité des liens à vérifier.\\\hline
CG05 - Développer ses apprentissages du mandat\par M0-C04, M0-C05 & Repérer un besoin d’apprentissage; chercher un appui; réexaminer une expérience. Buts : B2, B3. & E020 (S02, p. 87-88); E031 (S11, p. 12); E032 (S11, p. 12-13) Portée générale et nécessité des liens à vérifier.\\\hline
CG06 - Mobiliser les ressources d’un réseau syndical\par M0-C11 & Choisir un interlocuteur; formuler une demande d’appui; suivre l’échange. Buts : B1, B3. & E017 (S03, p. 217-218); E023 (S02, p. 90); E037 (S02, p. 106-107); E039 (S02, p. 110) Portée générale et nécessité des liens à vérifier.\\\hline
\end{longtable}\textbf{Frontières :} communiquer (CG02) soutient une tâche; préparer un accueil ou une instance vise un résultat particulier. Apprendre son mandat (CG05) se distingue de transmettre à une autre personne. Mobiliser un réseau (CG06) se distingue de coordonner les délégués dans une intervention locale.\par
\clearpage{\Large\bfseries\color{ink} Table de correspondance - compétences particulières}\par
\small Chaque ligne indique le parent M0, le résultat visé, les déterminants disponibles et les limites. Les indications ne constituent pas encore des objectifs de formation définitifs.\par
\renewcommand{\arraystretch}{1.12}\begin{longtable}{|p{61mm}|p{84mm}|p{115mm}|}\hline\rowcolor{pale}\textbf{Compétence et parent M0} & \textbf{Indication / résultat / buts} & \textbf{Déterminants et limites}\\\hline\endhead
CP-C01 - Préparer l’accueil des membres\par Parent : M0-C05 & Une présentation des ressources et du fonctionnement syndical. Buts : B1, B2. & E006 (S01, p. 307); E018 (S02, p. 80); E019 (S02, p. 81); E020 (S02, p. 87-88) Le volet information transversal est repris en CG02; l’accueil garde un résultat propre.\\\hline
CP-C02 - Réaliser une consultation des membres\par Parent : M0-C06 & Une synthèse des préoccupations et des avis reçus. Buts : B1, B2. & E008 (S01, p. 308-309); E021 (S02, p. 97-98); E027 (S04, p. 161-162) La représentativité de la consultation reste à examiner.\\\hline
CP-C03 - Préparer une instance syndicale\par Parent : M0-C02 & Une rencontre préparée avec des informations utiles à la décision. Buts : B1, B2. & E002 (S01, p. 301); E012 (S03, p. 212, 215, 219); E025 (S02, p. 102-103); E047 (S04, p. 145) La présidence de séance et les procédures ne sont pas entièrement décrites.\\\hline
CP-C04 - Traiter un dossier de représentation\par Parent : M0-C01 & Un dossier instruit et suivi avec les ressources concernées. Buts : B1. & E001 (S01, p. 301, 304); E005 (S01, p. 306-307); E016 (S03, p. 217); E030 (S04, p. 168) Le pouvoir de décision et les actes des spécialistes restent à distinguer.\\\hline
CP-C05 - Accompagner un membre en difficulté\par Parent : M0-C08 & Une aide syndicale délimitée et une orientation au besoin. Buts : B1, B2. & E009 (S01, p. 309-310); E015 (S03, p. 216-217) Le service rendu ne comprend pas un acte clinique.\\\hline
CP-C06 - Traiter un différend syndical\par Parent : M0-C07 & Une intervention ou un relais dont la raison est explicitée. Buts : B1. & E007 (S01, p. 308) Le périmètre de ce rôle est contesté dans les textes.\\\hline
CP-C07 - Coordonner le travail de l’équipe syndicale\par Parent : M0-C03 & Une répartition discutée des actions et des relais. Buts : B1, B2. & E003 (S01, p. 304); E013 (S03, p. 212, 215, 219); E014 (S03, p. 215-216, 219-220); E042 (S03, p. 229-232); E043 (S03, p. 230, 235) Coordonner ne suppose pas un pouvoir hiérarchique.\\\hline
CP-C08 - Organiser une action collective\par Parent : M0-C10 & Des moyens et des rôles préparés pour l’action. Buts : B1. & E011 (S01, p. 311); E022 (S02, p. 89) Le résultat de mobilisation dépend aussi du contexte.\\\hline
CP-C09 - Préparer une prise de position syndicale\par Parent : M0-C09 & Une position et un message préparés avec les acteurs concernés. Buts : B1. & E010 (S01, p. 310-311); E028 (S04, p. 162); E035 (S03, p. 221) La sensibilisation interne et le porte-parolat pourront être séparés.\\\hline
CP-C10 - Négocier les moyens de libération syndicale\par Parent : M0-C12 & Une demande de temps examinée avec les acteurs habilités. Buts : B1. & E024 (S02, p. 91); E045 (S03, p. 234-235) La demande, l’accord et la répartition sont trois éléments distincts.\\\hline
CP-C11 - Préparer une réponse à un projet de sous-traitance\par Parent : M0-C13 & Une réponse collective fondée sur les informations accessibles. Buts : B1. & E026 (S04, p. 159-161); E029 (S04, p. 167-168) Domaine spécialisé; pertinence pour le pilote à discuter.\\\hline
CP-C12 - Organiser une action de solidarité communautaire\par Parent : M0-C14 & Une action préparée avec une équipe et des partenaires. Buts : B1. & E038 (S02, p. 108) Exemple ponctuel, à examiner comme option plutôt que socle imposé.\\\hline
CP-C13 - Préparer la transmission d’un mandat syndical\par Parent : M0-C04 & Des connaissances et des relais préparés pour la relève. Buts : B2, B3. & E004 (S01, p. 305); E033 (S11, p. 13); E044 (S03, p. 231); E046 (S02, p. 125) L’apprentissage personnel est déplacé vers CG05; la transmission constitue ici le service rendu.\\\hline
\end{longtable}
\clearpage{\Large\bfseries\color{ink} Processus, portée et vérifications restantes}\par
\textbf{P1 - Recueillir les préoccupations des membres.} Repérer la demande ou le besoin à l’origine de l’intervention. Appuis partiels : E008 (S01, p. 308-309); E021 (S02, p. 97-98)\par
\textbf{P2 - Analyser les enjeux du mandat.} Examiner la situation et les limites du rôle exercé. Appuis partiels : E005 (S01, p. 306-307); E016 (S03, p. 217); E029 (S04, p. 167-168)\par
\textbf{P3 - Préparer l’intervention syndicale.} Préparer les informations, les moyens et les relais. Appuis partiels : E012 (S03, p. 212, 215, 219); E022 (S02, p. 89); E047 (S04, p. 145)\par
\textbf{P4 - Réaliser l’intervention syndicale.} Conduire l’action prévue dans le mandat concerné. Appuis partiels : E015 (S03, p. 216-217); E025 (S02, p. 102-103); E026 (S04, p. 159-161)\par
\textbf{P5 - Examiner les suites de l’intervention.} Reprendre les suites et les besoins de suivi; la séquence commune reste à établir. Appuis partiels : E001 (S01, p. 301, 304); E030 (S04, p. 168)\par
\textbf{Relations proposées :} les cercles suggèrent que l’activité générale contribue à la réalisation particulière; chaque cellule marquée doit encore être examinée. Le registre des relations du dépôt donne un motif pour chaque lien. Aucun lien ne repose sur une mesure d’efficacité ou un accord de terrain.\par
\textbf{Conditions distinctes des compétences :} cumul du métier et des responsabilités (E034), moyens de temps variables (E040), cumul des mandats (E041). Les conditions plus larges proposées par Gemini restent à documenter. Une difficulté liée aux moyens ne démontre pas un manque individuel de compétence.\par
\textbf{Fonctions :} les sources décrivent surtout des présidences, des groupes de fonctions ou des actions collectives. Aucune attribution automatique à la vice-présidence du SSPHQ. Les résultats S03/S11 peuvent se recouper.\par
\textbf{À vérifier avant un projet de formation arrêté :} pertinence de chaque compétence; résultat distinct; absence de regroupement artificiel; transferts des compétences générales; nécessité de chaque liaison; processus éventuel; seuil d’exercice attendu; apprentissages nécessaires; objectifs et durées; consultation et décision du responsable.\par
\textbf{Historique :} les parents M0 sont conservés. Les ajouts ou séparations sont tracés dans la table de correspondance. Les versions futures M1 et M2 pourront s’appuyer sur les données autorisées et les entretiens. Les états futurs restent vides dans le suivi.\par
\textbf{Relecture :} préparation documentaire et mise en forme par Codex le 26 octobre 2026. Examen personnel et choix de Steve à consigner. Cette présentation ne vaut pas conformité intégrale au guide ni validation ministérielle.
\end{document}
